\chapter{曲线积分与重积分}\label{ch:b2:multint}

本章将积分从区间推广到曲线，以及
平面与空间中的区域。曲线积分是对
\emph{微分形式} $P\,\dd x + Q\,\dd y$ 沿有向
弧的积分；二重与三重积分则分别对二维与
三维区域上的函数积分。两种理论在
\emph{格林--黎曼定理}，即二维的
微积分基本定理，全书的主要计算工具是
\emph{变量替换公式}，其伸缩因子为
雅可比行列式的绝对值。

\section{曲线积分}

\begin{definition}[微分形式；曲线积分]\label{def:b2:multint:lineint}
设 $U \subseteq \R^2$ 为开集。称\emph{微分
形式}\index{微分形式}为
$1$ 次、$\mathcal{C}^0$ 类、定义在 $U$ 上的
表达式 $\omega = P\,\dd x + Q\,\dd y$，其中 $P, Q \colon U \to \R$
连续 --- 形式上说，它是从 $U$ 到
$\R^2$ 的对偶空间的连续映射 $\omega(M) = P(M)\,e_1^* + Q(M)\,e_2^*$。对于
$\mathcal{C}^1$ 类的弧 $\gamma \colon [a, b] \to U$，$\gamma(t) =
(x(t), y(t))$，\emph{曲线积分}\index{曲线积分}
（即 $\omega$ 沿 $\gamma$ 的积分）
为
\[
\int_\gamma \omega
= \int_a^b \Bigl(P(\gamma(t))\,x'(t)
+ Q(\gamma(t))\,y'(t)\Bigr)\,\dd t .
\]
上述定义可原样推广到 $\R^3$（形式 $P\,\dd x +
Q\,\dd y + R\,\dd z$）及分段 $\mathcal{C}^1$ 弧（对
各段求和）。
\end{definition}

\begin{proposition}[不变性与定向]\label{prop:b2:multint:lineinv}
曲线积分在递增的 $\mathcal{C}^1$
参数变换下保持不变，在递减的参数变换下变号。
因此它只依赖于\emph{有向}几何弧。
\end{proposition}

\begin{proof}
若 $\theta \colon [c, d] \to [a, b]$ 为参数变换，且
$\tilde\gamma = \gamma \circ \theta$，则由链式法则及
一元积分的变量替换 $t = \theta(u)$，
\[
\int_{\tilde\gamma}\omega
= \int_c^d \bigl(P(\gamma(\theta(u)))\,x'(\theta(u))
+ Q(\gamma(\theta(u)))\,y'(\theta(u))\bigr)\,\theta'(u)\,\dd u
= \pm\int_a^b \bigl(Px' + Qy'\bigr)(t)\,\dd t ,
\]
符号为 $+$，当 $\theta$ 递增（$\theta(c) = a$）时；
为 $-$，递减时（积分上下限互换）。
\end{proof}

\begin{example}[力所做的功；环量]\label{ex:b2:multint:work}
若 $F = (P, Q)$ 为力场，$\int_\gamma P\dd x + Q\dd y =
\int_a^b \langle F(\gamma(t)), \gamma'(t)\rangle\,\dd t$ 是
$F$ 沿 $\gamma$ 所做的\emph{功}。对于 $\omega = -y\,\dd x +
x\,\dd y$ 沿逆时针单位圆 $\gamma(t) =
(\cos t, \sin t)$：
\[
\int_\gamma \omega
= \int_0^{2\pi}\bigl((-\sin t)(-\sin t)
+ \cos t\cos t\bigr)\,\dd t = 2\pi ,
\]
恰为所围面积的两倍 --- 格林--黎曼定理的初步体现。
\end{example}

\begin{example}[同一积分，两种参数化，一个
符号陷阱]\label{ex:b2:multint:semicircle}
计算 $\int_\gamma x\,\dd y$，沿上
单位半圆，从 $(1, 0)$ 到 $(-1, 0)$。取 $\gamma(t) =
(\cos t, \sin t)$，$t \in \intcc0\pi$：
\[
\int_0^\pi\cos t\cdot\cos t\,\dd t = \frac\pi2 .
\]
改用图像参数化 $x \mapsto (x, \sqrt{1 -
x^2})$，$x$ 从 $1$ 变到 $-1$（注意方向！）：
\[
\int_1^{-1}x\cdot\frac{-x}{\sqrt{1 - x^2}}\,\dd x
= \int_{-1}^{1}\frac{x^2}{\sqrt{1 - x^2}}\,\dd x
= \frac\pi2
\]
（令 $x = \sin u$，可化为沃利斯积分）。结果相同，
正如\cref{prop:b2:multint:lineinv}所保证的 --- 但仅因
两次积分都从 $(1,0)$ 到 $(-1,0)$；若改变行进方向，
则符号反转。沿 $x$ 轴补上闭合段
（此时 $\dd y = 0$）不产生贡献，而总和 $\frac\pi2$
恰为半圆盘的面积，这是下文格林--黎曼定理中
关于边界面积的公式的首个实例。
\end{example}

\begin{definition}[恰当形式与闭形式]\label{def:b2:multint:exact}
设 $\omega = P\,\dd x + Q\,\dd y$ 为 $\mathcal{C}^0$ 类的形式，
则称它在 $U$ 上为\emph{恰当形式}\index{恰当形式}，即存在 $f \in \mathcal{C}^1(U)$
（\emph{势函数}\index{势函数}）使 $\omega = \dd f$，即 $P = f_x$ 与 $Q
= f_y$。称 $\mathcal{C}^1$ 类的形式为\emph{闭形式}\index{闭形式}，若 $P_y = Q_x$ 在
$U$ 上成立。
\end{definition}

\begin{theorem}[曲线积分的
基本定理]\label{thm:b2:multint:ftc}
若 $\omega = \dd f$ 为恰当形式，$\gamma$ 是分段
$\mathcal{C}^1$ 弧，在 $U$ 中从 $A$ 到 $B$，则
\[
\int_\gamma \omega = f(B) - f(A) .
\]
特别地，恰当形式沿任一闭弧的积分
为零，且每个 $\mathcal{C}^1$ 恰当形式都是闭形式。
\end{theorem}

\begin{proof}
$\frac{\dd}{\dd t}f(\gamma(t)) = f_x(\gamma(t))x'(t) +
f_y(\gamma(t))y'(t)$，由链式法则
（\cref{ch:b2:diffcalc}）可得，故
\cref{def:b2:multint:lineint} 中的被积式即为 $t \mapsto
f(\gamma(t))$ 的导数，由微积分基本定理，
各段上结论即得；中间各项逐项相消。恰当
$\mathcal{C}^1$ 形式的闭性由施瓦茨定理给出：$P_y =
f_{xy} = f_{yx} = Q_x$。
\end{proof}

\begin{example}[重构势函数]\label{ex:b2:multint:potential}
设 $\omega = y\,\eu^{xy}\,\dd x + (x\,\eu^{xy} + 2y)\,\dd y$
在 $\R^2$ 上。它是闭形式：两个交叉偏导数都等于
$\eu^{xy}(1 + xy)$。求势函数时，先将 $P$ 关于
$x$ 积分，视 $y$ 为常数：
\[
f(x, y) = \int y\,\eu^{xy}\,\dd x = \eu^{xy} + c(y),
\]
再调整 $c$，比较 $f_y$：$x\,\eu^{xy} + c'(y) =
x\,\eu^{xy} + 2y$，得 $c(y) = y^2$。于是 $f(x,y) = \eu^{xy} +
y^2$。对于任一分段 $\mathcal C^1$ 弧，从 $(0,0)$
到 $(1,1)$，
\[
\int_\gamma\omega = f(1,1) - f(0,0) = (\eu + 1) - 1 = \eu ,
\]
与路径无关 --- 这一两步法（先对
$x$ 积分，再修正 $y$）是
\cref{thm:b2:multint:ftc} 在闭形式皆
恰当的区域上的实用逆命题。
\end{example}

\begin{example}[闭形式未必恰当]\label{ex:b2:multint:angleform}
在 $U = \R^2 \setminus \{0\}$ 上，称
\[
\omega = \frac{-y\,\dd x + x\,\dd y}{x^2 + y^2}
\]
为\emph{角形式}，它是闭形式（直接计算：$P_y$ 与 $Q_x$ 均等于
$\frac{y^2 - x^2}{(x^2+y^2)^2}$），但它沿单位
圆的积分为 $2\pi \neq 0$（与
\cref{ex:b2:multint:work} 的计算相同，除以 $1$）：$\omega$ 不是
恰当形式，在 $U$ 上。局部地，$\omega = \dd\theta$，其中
$\theta$ 为极角的一个分支；这一失效是整体性的 --- 角
无法在被挖去的点周围连续定义。对于
没有孔的区域，病态现象消失：在\emph{星形}
开集上，每个闭 $\mathcal{C}^1$ 形式都是恰当的
（庞加莱引理，见\cref{exo:b2:multint:8}）。
\end{example}

\section{二重积分}

一元黎曼积分的理论视为已知
（见本科一年级分册与\cref{ch:b2:integration}），并简述
其二元版本。设 $f$ 为矩形
$R = [a, b] \times [c, d]$ 上的连续函数，它有二重积分 $\iint_R f$，
其定义与一元情形完全一致，由网格上的黎曼和给出，并
可通过累次积分计算：

\begin{theorem}[矩形上的富比尼定理]\label{thm:b2:multint:fubini}
若 $f$ 在 $R = [a,b] \times [c,d]$ 上连续，
\[
\iint_R f
= \int_a^b \Bigl(\int_c^d f(x, y)\,\dd y\Bigr)\dd x
= \int_c^d \Bigl(\int_a^b f(x, y)\,\dd x\Bigr)\dd y .
\]
\end{theorem}

\begin{proof}
记 $F(x) = \int_c^d f(x, y)\,\dd y$。由 $f$ 在
紧集 $R$ 上一致连续，故 $F$ 连续（控制估计：
$\abs{F(x) - F(x')} \leq (d - c)\sup_y\abs{f(x,y) - f(x',y)}$）。
现将 $[a,b]$ 与 $[c,d]$ 各 $n$ 等分，得到一个
由小矩形 $R_{ij}$ 构成的网格，每个面积为 $\Delta x\,\Delta y$。对每个小矩形，有
$\inf_{R_{ij}} f \cdot \Delta x \Delta y \leq
\int_{x_{i-1}}^{x_i}\int_{y_{j-1}}^{y_j} f(x,y)\,\dd y\,\dd x \leq
\sup_{R_{ij}} f \cdot \Delta x \Delta y$，由一元积分的
单调性（两次应用）。对所有小矩形求和，
累次积分 $\int_a^b F$ 夹在网格的下
黎曼和与上黎曼和之间；由一致连续性，两个和
都收敛于定义 $\iint_R f$ 的共同值，当 $n \to
\infty$ 时。将 $x$ 与 $y$ 的角色互换，
同一论证同样适用，故两个累次积分都等于 $\iint_R f$。
\end{proof}

\begin{remark}
紧矩形上的连续性在
富比尼定理的证明中起实质作用：它提供了
夹住黎曼和的一致连续性。对更不规则
的被积函数，结论确实不成立 --- 存在这样的函数，其两个
累次积分都存在而不相等。真正的一般
定理仅以可积性为唯一假设，即
勒贝格积分下的富比尼定理，
在本科三年级分册中证明；本章全部内容始终
限于连续情形，在此情形下
上述初等证明已完整。
\end{remark}

\begin{definition}[初等区域]\label{def:b2:multint:domain}
称区域 $D \subseteq \R^2$ 为\emph{$y$-初等的}，若
\[
D = \{(x, y) : a \leq x \leq b,\
\varphi_1(x) \leq y \leq \varphi_2(x)\}
\]
其中 $\varphi_1 \leq \varphi_2$ 在 $[a,b]$ 上连续
（$x$-初等的情形对称地定义）。设 $f$ 在
$y$-初等区域 $D$ 上连续，则
\[
\iint_D f
= \int_a^b\Bigl(
\int_{\varphi_1(x)}^{\varphi_2(x)} f(x,y)\,\dd y\Bigr)\dd x ,
\]
可以验证（通过逼近延拓 $f$，或
通过分割）：当 $D$ 在两个方向上均为初等时，两个
累次积分相等。分割为有限个初等部分的
区域由可加性处理，并
定义\emph{面积}\index{面积!平面区域的}：$D$ 的面积为 $\operatorname
{Area}(D) = \iint_D 1$。
\end{definition}

\begin{example}\label{ex:b2:multint:triangle}
在三角形 $D = \{0 \leq x \leq 1,\ 0 \leq y \leq x\}$ 上：
\[
\iint_D xy \,\dd x\,\dd y
= \int_0^1 x\Bigl(\int_0^x y\,\dd y\Bigr)\dd x
= \int_0^1 x\cdot\frac{x^2}{2}\,\dd x = \frac18 .
\]
交换积分次序（$x$ 从 $y$ 变到 $1$）：$\int_0^1
y\bigl(\int_y^1 x\,\dd x\bigr)\dd y = \int_0^1
y\,\frac{1 - y^2}{2}\,\dd y = \frac18$ --- 结果相同，
但计算方式不同：选取好的积分次序，
便是此中技艺的一半。
\end{example}

\begin{example}[只有一种次序可行]\label{ex:b2:multint:swaporder}
计算 $I = \displaystyle\int_0^1\!\!\int_x^1
\eu^{y^2}\,\dd y\,\dd x$。按原式所写，内层积分
$\int\eu^{y^2}\dd y$ 没有初等原函数，
计算便无法继续。但积分区域是三角形 $0 \leq
x \leq y \leq 1$，两个方向上均为初等区域；
交换积分次序，
\[
I = \int_0^1\!\!\int_0^y \eu^{y^2}\,\dd x\,\dd y
= \int_0^1 y\,\eu^{y^2}\,\dd y
= \Bigl[\tfrac12\eu^{y^2}\Bigr]_0^1 = \frac{\eu - 1}{2} .
\]
内层变量 $x$ 在被积函数中完全不出现，因此
先对它积分便恰好得到因子 $y$，
使外层积分立即可算。由此可见，富比尼定理不仅是
允许累次计算的许可 --- 更是允许\emph{选择}的许可，
正确的次序可以把原本算不出的积分化为
一行即可算完的式子。开始之前，务必先画出区域草图，并读出两种
描述方式。
\end{example}

\begin{theorem}[变量替换]\label{thm:b2:multint:changevar}
设 $\Phi \colon U' \to U$ 为 $\mathcal{C}^1$ 微分同胚，
它是 $\R^2$ 中开集之间的映射，设 $K \subseteq U$ 为紧区域，
可分割为初等部分，且 $K' = \Phi^{-1}(K)$，
再设 $f$ 在 $K$ 上连续，则
\[
\iint_K f(x, y)\,\dd x\,\dd y
= \iint_{K'} f\bigl(\Phi(u, v)\bigr)\,
\abs{\det J_\Phi(u, v)}\,\dd u\,\dd v .
\]
\end{theorem}
\admitted

\begin{remark}
完整证明 --- 在细网格上用微分逼近 $\Phi$，
并控制边界小格 --- 虽然冗长，
但并不深刻；它在本科三年级的测度论中完整给出，
作为勒贝格理论的推论。启发性的图像，正是
曲面面积中已用过的那幅：边长为 $\dd u$ 的小正方形位于
$(u, v)$ 处，一阶近似下被映为由 $\Phi_u\,\dd u$ 与 $\Phi_v\,\dd v$ 张成的平行四边形，
其面积为
$\abs{\det J_\Phi}\,\dd u\,\dd v$
（见\cref{lem:b2:surfaces:lagrange}）。
\end{remark}

\begin{remark}[方法：选取变量替换]
三种思路覆盖大多数情形。\emph{被积函数的
对称性}：$x^2 + y^2$ 宜用极坐标，乘积结构
则宜保留直角坐标轴。\emph{边界的
形状}：边界 $u(x,y) = c_1$、$v(x,y) = c_2$ 宜直接取
坐标 $(u, v)$ 本身，正如下文的双曲区域例子 ---
区域变为矩形，
这正是全部的胜利。\emph{线性结构}：
含 $x + y$ 与 $x - y$ 的式子，宜用 $45$ 度的
旋转或错切（\cref{ex:b2:multint:affine}）。
一切情形下，积分前须核对三点：映射
是新区域到旧区域上的双射；其雅可比
行列式须按\emph{实际使用的方向}计算（若反向更易，
则最后求逆）；且雅可比行列式以其
绝对值进入计算。
\end{remark}

\begin{example}[极坐标]\label{ex:b2:multint:polar}
$\Phi(\rho, \alpha) = (\rho\cos\alpha,\ \rho\sin\alpha)$ 的雅可比矩阵为
\[
J_\Phi = \begin{pmatrix}
\cos\alpha & -\rho\sin\alpha\\
\sin\alpha & \rho\cos\alpha
\end{pmatrix},
\qquad \det J_\Phi = \rho ,
\]
于是 $\dd x\,\dd y = \rho\,\dd\rho\,\dd\alpha$。对圆盘 $D_R$
（半径为 $R$）：
\[
\iint_{D_R} e^{-(x^2 + y^2)}\,\dd x\,\dd y
= \int_0^{2\pi}\!\!\int_0^R e^{-\rho^2}\rho\,\dd\rho\,\dd\alpha
= \pi\bigl(1 - e^{-R^2}\bigr)
\xrightarrow[R\to\infty]{} \pi .
\]
与正方形 $[-R, R]^2$ 比较（它夹在
圆盘 $D_R$ 与 $D_{R\sqrt2}$ 之间，被积函数皆为正）得
$\bigl(\int_{-\infty}^\infty e^{-x^2}\dd x\bigr)^2 = \pi$：
\[
\boxed{\ \int_{-\infty}^{+\infty} e^{-x^2}\,\dd x = \sqrt{\pi}\ }
\]
--- 这又是高斯积分，此次采用其最著名的证明
（比较单变量情形下的推导，见
\cref{ch:b2:integration}）。
\end{example}

\begin{example}[仿射变量替换]\label{ex:b2:multint:affine}
对仿射映射 $\Phi(u, v) = M(u, v)^{\mathsf T} + C$，其中
$M$ 可逆，雅可比矩阵为常矩阵 $M$：
面积乘以常数因子 $\abs{\det M}$
--- 在\cref{ch:b2:affine}中所作的承诺，如今已成为定理。
两个直接的应用。椭圆 $\frac{x^2}{a^2} +
\frac{y^2}{b^2} \leq 1$ 为单位圆盘在
$(u, v) \mapsto (au, bv)$ 下的像，故其面积为 $ab \cdot \pi$ ---
无需计算。对于 $f(x + y)$ 在
正方形 $K = \intcc01^2$ 上的积分，错切 $\Phi(u, v) = (u - v, v)$
（行列式为 $1$）将其化为 $f(u)$ 在一个
平行四边形上的积分，由富比尼定理沿常数 $u$ 切片计算。取 $f =
\exp$，
\[
\iint_K \eu^{x+y}\,\dd x\,\dd y
= \Bigl(\int_0^1 \eu^x\,\dd x\Bigr)^2 = (\eu - 1)^2,
\]
如乘积结构所确认的那样。选取
适合被积函数 --- 而非区域 --- 的坐标，
是此中技艺的另一半。
\end{example}


\begin{example}[适应曲线型区域的
坐标]\label{ex:b2:multint:hyperbolicregion}
设 $D$ 为第一象限中由
双曲线 $xy = 1$ 与 $xy = 3$ 以及直线 $y = x$ 与
$y = 3x$ 所围成。在坐标 $u = xy$、$v = y/x$ 下，区域
变为正方形 $\intcc13 \times \intcc13$；反解得
\[
x = \sqrt{u/v}, \qquad y = \sqrt{uv},
\qquad
\det J = x_uy_v - x_vy_u = \frac{1}{2v}
\]
（两行计算，取 $x = u^{1/2}v^{-1/2}$，$y =
u^{1/2}v^{1/2}$）。由此
\[
\operatorname{Area}(D)
= \int_1^3\!\!\int_1^3\frac{\dd u\,\dd v}{2v}
= 2\cdot\frac{\ln 3}{2} = \ln 3 \approx 1.10 .
\]
在直角坐标中切分 $D$，就意味着把它切成
三块，其边界含双曲线与直线 ---
可行、枯燥且易出错。这一道理与
\cref{ex:b2:multint:affine} 所述完全相同，且更为充分：读出
边界方程，让\emph{它们}来选取
坐标；雅可比行列式随即将曲线网格的
单元面积换算过来，恰如 $\rho$ 在
极坐标中所做的那样。
\end{example}

\begin{example}[平均值]\label{ex:b2:multint:average}
$f$ 在区域 $D$ 上的\emph{平均值}是
$\frac1{\operatorname{Area}(D)}\iint_Df$。例如，
均匀取点于圆盘内时，该点到圆心的平均距离
（圆盘半径为 $R$）为
\[
\frac{1}{\pi R^2}\int_0^{2\pi}\!\!\int_0^R
\rho\cdot\rho\,\dd\rho\,\dd\alpha
= \frac{2\pi R^3/3}{\pi R^2} = \frac{2R}3 ,
\]
不是 $R/2$：均匀面积使大半径处占有更多的
比重（半径为 $\rho$ 的圆环比重与 $\rho$ 成正比），
故平均值超过半径的一半。要算对这一
系数，恰是极坐标雅可比行列式在起作用，而
同样的权重也解释了本章后文中半球的形心 $\bar z = 3R/8$，
而非 $R/2$。
\end{example}

\section{格林--黎曼定理}

\begin{theorem}[格林--黎曼]\label{thm:b2:multint:green}
设 $K \subseteq \R^2$ 为紧区域，它在
两个方向上都是初等的（或若干此类区域沿线段粘合而成的有限并），
边界 $\partial K$ 为分段 $\mathcal{C}^1$ 的闭曲线，
并按\emph{逆时针}方向定向（区域位于
左侧）。设 $P, Q$ 为 $\mathcal{C}^1$ 类函数，
且在 $K$ 的某邻域上有定义，则：
\[
\oint_{\partial K} P\,\dd x + Q\,\dd y
= \iint_K \Bigl(\frac{\partial Q}{\partial x}
- \frac{\partial P}{\partial y}\Bigr)\,\dd x\,\dd y .
\]
\end{theorem}

\begin{proof}
首先，沿线段将 $K$
切成两块 $K_1, K_2$，等式两边都具有可加性：
二重积分由 $\iint$ 的可加性相加；至于边界上的积分，
$K_1$、$K_2$ 的逆时针边界各
经过内部切口一次，方向\emph{相反}，因此在
和式中
\[
\oint_{\partial K_1} + \oint_{\partial K_2}
= \oint_{\partial K} + (\text{割线，往返各一次})
= \oint_{\partial K},
\]
沿切口的两次经过互相抵消
（见\cref{prop:b2:multint:lineinv}），仅外边界
保留。经有限次切割，只需讨论
初等区域的情形。证明 $\oint P\,\dd x = -\iint_K P_y$ 在
$y$ 型初等区域 $D = \{a \leq x \leq b,\ \varphi_1(x) \leq y
\leq \varphi_2(x)\}$ 上成立；恒等式 $\oint Q\,\dd y = \iint_K Q_x$
同理可证（$x$ 型初等），定理即二者之和。

用富比尼定理及一元函数的
微积分基本定理计算二重积分：
\[
\iint_D \frac{\partial P}{\partial y}\,\dd x\,\dd y
= \int_a^b \bigl(P(x, \varphi_2(x)) - P(x, \varphi_1(x))\bigr)
\,\dd x .
\]
现在 $D$ 的逆时针边界由以下部分组成：下方的
图像 $y = \varphi_1(x)$，自左向右走；右侧的
竖直线段 $x = b$（向上），上方的图像 $y =
\varphi_2(x)$，\emph{自右向左}走过；左侧的竖直
线段 $x = a$（向下）。沿竖直线段时 $x$
为常数，故竖直段的贡献为 $0$；对 $\oint P\,\dd x$ 而言，图像
以 $x$ 为参数，给出
\[
\oint_{\partial D} P\,\dd x
= \int_a^b P(x, \varphi_1(x))\,\dd x
- \int_a^b P(x, \varphi_2(x))\,\dd x
= -\iint_D \frac{\partial P}{\partial y}\,\dd x\,\dd y .
\qedhere
\]
\end{proof}

\begin{corollary}[由边界求面积]\label{cor:b2:multint:area}
在\cref{thm:b2:multint:green}的条件下，
\[
\operatorname{Area}(K)
= \oint_{\partial K} x\,\dd y
= -\oint_{\partial K} y\,\dd x
= \frac12\oint_{\partial K} x\,\dd y - y\,\dd x .
\]
\end{corollary}

\begin{proof}
分别取 $(P, Q) = (0, x)$、$(-y, 0)$ 及
$\frac12(-y, x)$，应用格林–黎曼公式，每次都有 $Q_x - P_y = 1$。
\end{proof}

\begin{remark}[三种面积公式的选取]
三个边界公式的值相等，但在实际应用中
并不能互相替代。用 $\oint x\,\dd y$，当参数化使 $\dd y$ 简单时
（如以 $y$ 为自变量的图像）；对称地，
用 $-\oint y\,\dd x$；当参数化对 $x$ 与 $y$ 同等对待时，
则取对称的半和。以椭圆为例，
它给出一个常值被积函数，全然无需三角
线性化。在多边形边界上，半和
化为\cref{exo:b2:multint:12}中的鞋带公式，
即测量员的算法。当给定的参数化使边界
按顺时针方向走过时，三个公式
都给出面积的\emph{相反数}：负值结果并非
计算错误，而是定向的提示——应改变
符号，或改换参数化。
\end{remark}

\begin{example}[椭圆的面积]\label{ex:b2:multint:ellipsearea}
取 $x = a\cos t$，$y = b\sin t$，$t \in [0, 2\pi]$：
\[
\operatorname{Area}
= \frac12\int_0^{2\pi}\bigl(a\cos t \cdot b\cos t
- b\sin t\cdot(-a\sin t)\bigr)\,\dd t
= \frac{ab}{2}\int_0^{2\pi}\dd t = \pi ab .
\]
\end{example}

\begin{example}[用格林–黎曼公式作
验算]\label{ex:b2:multint:greencheck}
取 $P = -y^3$、$Q = x^3$，定义在闭单位圆盘 $D$ 上。
边界一侧，取 $\gamma(t) = (\cos t, \sin t)$：
\[
\oint_{\partial D}P\,\dd x + Q\,\dd y
= \int_0^{2\pi}\bigl(\sin^4 t + \cos^4 t\bigr)\dd t
= 2\pi\cdot\Bigl(\frac38 + \frac38\Bigr) = \frac{3\pi}2 ,
\]
由线性化（$\sin^4 + \cos^4 = \tfrac34 +
\tfrac14\cos4t$）。内部一侧：
\[
\iint_D(Q_x - P_y)\,\dd x\,\dd y
= \iint_D 3(x^2 + y^2)\,\dd x\,\dd y
= 3\int_0^{2\pi}\!\!\int_0^1\rho^3\,\dd\rho\,\dd\alpha
= \frac{3\pi}2 .
\]
结果相同，计算却截然不同——这正是
其实用之处：格林公式的两侧中，哪一侧更易，
便作为计算，另一侧作为验证。对于多项式场沿
闭曲线的环量，二重积分几乎总是
更容易的一侧。
\end{example}

\begin{remark}
格林–黎曼公式解释了\cref{ex:b2:multint:angleform}：
对闭形式（$Q_x = P_y$）而言，任一包含在 $U$ 中的区域
之边界上的积分为零。角形式不是恰当形式，只因
原点处的穿孔使以单位圆为边界的圆盘无法
位于 $U$ 之内——闭形式的线积分
可探测区域的孔洞。（若再深入，
此观察即发展为德拉姆上同调。）
\end{remark}

\section{三重积分}

理论推广到三个变量并无新思想：富比尼定理
将 $\iiint$ 化为三次一元积分（可
用\emph{截面法}：$\iiint_K f = \int\bigl(\iint_{K_z}
f\bigr)\dd z$，沿水平截面 $K_z$ 进行，或
用\emph{堆叠法}：沿竖直细条先对 $z$ 积分），
并且变量替换公式在 $3 \times 3$
雅可比矩阵下仍然成立。

\begin{example}[柱坐标与球坐标]\label{ex:b2:multint:spherical}
\emph{柱坐标}：$(x, y, z) = (\rho\cos\alpha, \rho\sin\alpha,
z)$，$\dd x\,\dd y\,\dd z = \rho\,\dd\rho\,\dd\alpha\,\dd z$。
\emph{球坐标}：$(x, y, z) = (r\cos\theta\cos\varphi,\
r\sin\theta\cos\varphi,\ r\sin\varphi)$（$\theta$ 为经度，
$\varphi \in [-\frac\pi2, \frac\pi2]$ 为纬度）：按最后一行展开 $3
\times 3$ 行列式，
\[
\det J = r^2\cos\varphi ,
\qquad
\dd x\,\dd y\,\dd z
= r^2\cos\varphi\;\dd r\,\dd\theta\,\dd\varphi .
\]
半径为 $R$ 的球的体积：
\[
V = \int_0^R\!\!\int_0^{2\pi}\!\!\int_{-\pi/2}^{\pi/2}
r^2\cos\varphi\;\dd\varphi\,\dd\theta\,\dd r
= \frac{R^3}{3}\cdot 2\pi \cdot 2
= \boxed{\frac43\pi R^3} ,
\]
终于兑现了前几册体积章节中
所承认的公式。
\end{example}

\begin{example}[四面体的两次计算]
\label{ex:b2:multint:tetra}
求 $T = \{x, y, z \geq 0,\ x + y + z \leq 1\}$ 的体积，
用堆叠法：对固定的 $(x, y)$（位于三角形 $x + y \leq 1$ 内），
$z$ 取遍 $\intcc0{1 - x - y}$，于是
\[
V = \int_0^1\!\!\int_0^{1-x}(1 - x - y)\,\dd y\,\dd x
= \int_0^1\frac{(1 - x)^2}{2}\,\dd x = \frac16 .
\]
用截面法：高度为 $z$ 的截面是三角形 $\{x, y
\geq 0,\ x + y \leq 1 - z\}$，面积为 $\frac{(1-z)^2}2$，且
$V = \int_0^1\frac{(1-z)^2}2\,\dd z = \frac16$，结果相同——两种
计算仅是同一积分的次序不同，
这正是富比尼定理所断言的全部。$\frac16 =
\frac13\cdot\frac12\cdot1$ 恰为底面为三角形的锥体
公式（见\cref{ex:b2:multint:cone}）；$n$ 维推广 $1/n!$
也正是由同样的截面法证明，
在周末问题中给出。
\end{example}

\begin{example}[半球体的形心]\label{ex:b2:multint:centroid}
设 $H$ 为半径为 $R$ 的上半球体（$z \geq 0$），
其形心高度为 $\bar z = \frac1{V}\iiint_H z$，其中 $V =
\frac23\pi R^3$。采用球坐标（$z =
r\sin\varphi$，$\varphi \in \intcc0{\pi/2}$）：
\[
\iiint_H z
= \int_0^R r^3\,\dd r\int_0^{2\pi}\dd\theta
\int_0^{\pi/2}\sin\varphi\cos\varphi\,\dd\varphi
= \frac{R^4}4\cdot2\pi\cdot\frac12 = \frac{\pi R^4}4 ,
\]
于是
\[
\bar z = \frac{\pi R^4/4}{2\pi R^3/3} = \frac{3R}8 :
\]
实心半球的形心位于平底面以上
半径的八分之三处——低于半高
$R/2$，正如所料，因为该立体靠近底面处更厚实。
每一次形心计算都具有这样的结构：一个力矩
积分、一个体积、一个比值，以及一项合理性检验，
对照几何直观加以核验。
\end{example}

\begin{example}[用截面法求体积：锥体]\label{ex:b2:multint:cone}
设锥体底面积为 $A$，高为 $h$（顶点朝上，底面位于 $z = 0$）：
高度为 $z$ 处的截面，是将底面按因子 $(1 -
z/h)$ 缩放所得，面积为 $A(1 - z/h)^2$。由此
\[
V = \int_0^h A\Bigl(1 - \frac zh\Bigr)^2\dd z = \frac{Ah}{3} :
\]
这是中学课本中的“三分之一”公式，对\emph{任意}底面形状均成立
——截面法将其化为一个平方的积分。
\end{example}

\begin{example}[平面上的
可积性阈值]\label{ex:b2:multint:threshold}
对哪些 $\alpha > 0$，$\iint_{D}\rho^{-\alpha}\,\dd
x\,\dd y$ 在去心单位圆盘 $D$ 上收敛（极限
沿环域 $\varepsilon \leq \rho \leq 1$ 取值）？
在极坐标下，
\[
\int_0^{2\pi}\!\!\int_\varepsilon^1\rho^{-\alpha}\,
\rho\,\dd\rho\,\dd\alpha
= 2\pi\int_\varepsilon^1\rho^{1-\alpha}\,\dd\rho ,
\]
当 $\varepsilon \to 0$ 时，该积分收敛当且仅当 $1 - \alpha > -1$，
即 $\alpha < 2$：在 $2$ 维中，临界奇异指数恰为空间维数本身，
雅可比行列式带来的额外因子 $\rho$
使奇异性减弱一个幂次（类似地，
对空间中的点奇点有 $\alpha < 3$，借助 $r^2$ 即可得出。）
这类径向估算正是判定可积性的办法，
在本科三年级的勒贝格框架下一眼即可看出——
这也正是 $\iiint 1/r$ 能
在\cref{exo:b2:multint:7}中轻松收敛的原因。
\end{example}

\begin{remark}[常见易错点]
(i) \emph{定向}：曲线积分
随行进方向改变符号，格林--黎曼公式要求边界
取逆时针方向（区域位于左侧）；对于带洞的区域，
其内边界须取\emph{顺时针}方向。(ii)
\emph{雅可比行列式取绝对值}：换元
永远不会产生负的面积，而漏掉
$\abs{\det}$ 通常恰好在映射
反转定向时使符号翻转。(iii) \emph{极坐标因子 $\rho$}：
$\dd x\,\dd y = \rho\,\dd\rho\,\dd\alpha$，而不是
$\dd\rho\,\dd\alpha$——这是全章最常见的错误，
量纲分析可以识别它，因为
$\dd\rho\,\dd\alpha$ 的量纲是长度而非
面积。(iv) \emph{反常二重积分}：取极限时
沿增长的圆盘与增长的方形所得极限一致，因为
被积函数恒正（夹逼）；对于变号的被积函数，
极限可能依赖于穷竭方式，因此
在无绝对收敛的前提下不作任何断言。(v)
\emph{区域与被积函数}：乘积形式的被积函数在
非乘积区域上并\emph{不}能分解积分——
分解需要二者同时满足，如\cref{ex:b2:multint:polar}
中的正方形所示。
\end{remark}

\begin{remark}[本册内的几个视角]
此处计算的高斯积分悄然遍布于
概率各章：常数 $\sqrt\pi$ 藏于
斯特林公式（\cref{thm:b2:comparison:stirling}）中，它正是
本章的这一积分，经由斯特林公式，决定了
$1/\sqrt{\pi n}$ 的渐近行为，即随机游走
在\cref{ch:b2:proba}中的返回概率。周末问题中的沃利斯积分
同样在那里再次出现，
驱动相同的中心二项式估计。反过来，本章的面积元
与体积元完善了
\cref{ch:b2:surfaces}中的几何，而格林公式重新计算了
\cref{ch:b2:curves}中的包络面积（星形线，
见\cref{exo:b2:multint:5}）。一章三用：
为几何提供度量，为概率提供常数，以及
二者共同依赖的换元规范。
\end{remark}

\section{练习}

\begin{exercise}[$\star$]\label{exo:b2:multint:1}
计算 $\int_\gamma y^2\,\dd x + x\,\dd y$，沿以下路径：(a) 线段
从 $(0,0)$ 到 $(1,1)$；(b) 抛物线弧 $y = x^2$ 从
$(0,0)$ 到 $(1,1)$。该微分形式是否为恰当形式？
\end{exercise}

\begin{exercise}[$\star$]\label{exo:b2:multint:2}
证明 $\omega = (2xy + y^3)\,\dd x + (x^2 + 3xy^2 + 1)\,\dd y$
在 $\R^2$ 上是闭形式，求其势函数，并计算
$\int_\gamma\omega$，沿任一从 $(0, 0)$ 到 $(1, 2)$ 的几何弧。
\end{exercise}

\begin{exercise}[$\star$]\label{exo:b2:multint:3}
计算 $\iint_D (x + y)\,\dd x\,\dd y$，其中 $D$ 是
由 $y = x^2$ 与 $y = x$ 所围成的区域（$0 \leq x \leq 1$），
分别按两种积分次序计算。
\end{exercise}

\begin{exercise}[$\star\star$]\label{exo:b2:multint:4}
利用极坐标，计算 $\iint_D \frac{\dd x\,\dd y}{(1 +
x^2 + y^2)^2}$ 在整个平面上的积分（视为对圆盘的极限），以及
$\iint_{D'} xy\,\dd x\,\dd y$ 在四分之一圆盘 $D' = \{x, y
\geq 0,\ x^2 + y^2 \leq 1\}$ 上的积分。
\end{exercise}

\begin{exercise}[$\star\star$]\label{exo:b2:multint:5}
计算星形线 $x = \cos^3 t$，$y =
\sin^3 t$，$t \in [0, 2\pi]$ 所围图形的面积，利用
\cref{cor:b2:multint:area}。\emph{（将 $\sin^2 t\cos^2 t$ 线性化。）}
\end{exercise}

\begin{exercise}[$\star\star$]\label{exo:b2:multint:6}
计算由下方的抛物面 $z
= x^2 + y^2$ 与上方的平面 $z = 1$ 所围立体的体积，分别用两种方法计算：
叠加法（对 $1 - x^2 - y^2$ 在单位圆盘上积分，用
极坐标）与切片法（水平截面为半径
$\sqrt z$）。
\end{exercise}

\begin{exercise}[$\star\star$]\label{exo:b2:multint:7}
（球体的引力 --- 牛顿定理的特殊情形
）证明球壳 $a \leq r \leq
b$ 的体积为 $\frac43\pi(b^3 - a^3)$，并计算 $\iiint_{B}
\frac{\dd x\,\dd y\,\dd z}{r}$ 在球 $B$（半径为 $R$）上的积分
（$r$ 为到原点的距离）。\emph{（
球坐标。）}
\end{exercise}

\begin{exercise}[$\star\star\star$]\label{exo:b2:multint:8}
（庞加莱引理，星形情形）设 $U$ 为
关于 $0$ 的星形区域（即 $M \in U \Rightarrow [0, M] \subseteq U$），
且 $\omega = P\dd x + Q\dd y$ 为闭的 $\mathcal{C}^1$ 形式，
定义于 $U$ 上。令
\[
f(x, y) = \int_0^1 \bigl(x\,P(tx, ty) + y\,Q(tx, ty)\bigr)\dd t .
\]
利用积分号下的求导
（\cref{ch:b2:integration}）及 $P_y = Q_x$，证明 $f_x = P$
与 $f_y = Q$：星形开集上的每个闭形式
都是恰当形式。
\end{exercise}

\begin{exercise}[$\star\star\star$]\label{exo:b2:multint:9}
（通过二重积分计算狄利克雷积分）论证并利用
\[
\int_0^\infty\!\!\int_0^\infty e^{-xy}\sin x\;\dd y\,\dd x
\quad\text{与}\quad
\int_0^\infty\!\!\int_0^\infty e^{-xy}\sin x\;\dd x\,\dd y
\]
在 $[0, A] \times [0, \infty)$ 上，证明 $\int_0^A \frac{\sin
x}{x}\dd x = \frac\pi2 - \int_0^\infty
e^{-Ay}\frac{y\sin A + \cos A}{1 + y^2}\dd y$，并由此推出
$\int_0^\infty \frac{\sin x}{x}\,\dd x = \frac\pi2$，并与
含参积分的证明作比较（参见
\cref{ch:b2:integration}）。
\end{exercise}

\begin{exercise}[$\star\star\star$]\label{exo:b2:multint:10}
（等周不等式的维廷格证明）设 $\gamma$ 为简单
闭 $\mathcal{C}^1$ 曲线，长度为 $2\pi$，以弧长为参数，
定义在 $[0, 2\pi]$ 上，所围面积为 $A$。利用
\cref{cor:b2:multint:area}、帕塞瓦尔恒等式与维廷格不等式，
（\cref{ch:b2:fourier}中的练习），证明 $A \leq \pi$，
并在圆周时取等号。\emph{（归一化，使 $\int_0^{2\pi} x(s)\dd s
= 0$；写出 $2A = \oint x\,\dd y - y\,\dd x$，并估计 $2A \leq
\int (x^2 + y'^2)$，细致地借助 $2A = \int_0^{2\pi}(xy' - yx')\dd
s$ 及 $x^2 + y'^2 \geq 2xy'$。）}
\end{exercise}

\begin{exercise}[$\star\star$]\label{exo:b2:multint:11}
（球的矩）对球 $B$（半径为 $R$，在
$\R^3$ 中），在球坐标下计算 $\iiint_B z^2\,\dd x\,\dd y\,\dd z$，
并由对称性推出 $\iiint_B (x^2 + y^2 +
z^2)\,\dd x\,\dd y\,\dd z$ 的值。将后者与
球壳计算 $\int_0^R r^2\cdot4\pi
r^2\,\dd r$ 作核对。
\end{exercise}

\begin{exercise}[$\star\star$]\label{exo:b2:multint:12}
（鞋带公式）设 $K$ 为多边形，顶点为 $(x_1,
y_1), \dots, (x_m, y_m)$，按逆时针顺序排列（下标
模 $m$）。由\cref{cor:b2:multint:area}可得：
\[
\operatorname{Area}(K)
= \frac12\sum_{i=1}^m
\bigl(x_iy_{i+1} - x_{i+1}y_i\bigr) ,
\]
并以 $(0,0)$、$(1,0)$、
$(0,1)$ 为顶点的三角形检验该公式。
\end{exercise}

\section{问题：$n$ 维球的体积}

\begin{problem}[{周末问题 --- $V_n =
\pi^{n/2}/\Gamma(\frac n2 + 1)$，以及高维的
奇异现象}]\label{pb:b2:multint:1}
圆盘的面积为 $\pi$，球的体积为 $\frac43\pi$ --- 那么
呢？本题计算
$\R^n$ 中单位球的体积，对每个 $n$ 各算两次（一次用由
沃利斯积分驱动的切片递推，另一次借助 $\Gamma$ 函数及
\cref{ex:b2:multint:polar}中的高斯积分），然后读出
几何结论：体积在五维时达到峰值，随后迅速
趋于零，而高维球的几乎全部体积都藏在
靠近边界的薄壳中。\index{球体积}\index{沃利斯
积分}对 $\R^n$ 中某球上的连续函数，
其积分理解为 $n$ 重累次积分
（每次只对一个坐标做切片，与 $n
\leq 3$ 时的做法一致）。记 $B_n(R)$ 为半径为 $R$ 的闭球，
中心在 $0$，其体积为 $v_n(R)$，并记 $V_n =
v_n(1)$，约定 $V_0 = 1$。

\medskip
\noindent\textbf{第一部分 --- 切片递推。}
\begin{enumerate}
  \item 将 $x_i = Ru_i$ 代入 $n$ 个累次
        积分中的每一个，证明 $v_n(R) = V_nR^n$。
  \item 沿最后一个坐标对 $B_n(1)$ 作切片，证明
        \[
        V_n = V_{n-1}\int_{-1}^{1}(1 -
        t^2)^{\frac{n-1}2}\,\dd t .
        \]
  \item 令 $t = \sin\theta$，将该积分识别为
        沃利斯积分：$\int_{-1}^1(1 -
        t^2)^{\frac{n-1}2}\dd t = 2W_n$，其中 $W_n =
        \int_0^{\pi/2}\cos^n\theta\,\dd\theta =
        \int_0^{\pi/2}\sin^n\theta\,\dd\theta$。
  \item 证明两个沃利斯恒等式（用分部积分法；
        再对 $nW_nW_{n-1}$ 作裂项相消）：
        \[
        W_n = \frac{n-1}nW_{n-2} \quad (n \geq 2),
        \qquad
        W_nW_{n-1} = \frac{\pi}{2n} \quad (n \geq 1).
        \]
\end{enumerate}

\medskip
\noindent\textbf{第二部分 --- 递推式的求解。}
\begin{enumerate}[resume]
  \item 综合第 2--4 题，得到两步递推式
        \[
        V_n = \frac{2\pi}{n}\,V_{n-2}
        \qquad (n \geq 2).
        \]
  \item 对 $k \geq 0$，推出以下闭式：
        \[
        V_{2k} = \frac{\pi^k}{k!},
        \qquad
        V_{2k+1} = \frac{2^{k+1}\pi^k}{1\cdot3\cdot5\cdots
        (2k+1)} .
        \]
  \item 数值列出 $V_1, \dots, V_7$ 的值。利用
        比值 $V_n/V_{n-2} = 2\pi/n$ 及 $2W_5$
        与 $2W_6$ 的值，证明数列 $(V_n)$ 递增，
        直至其最大值 $V_5 = \frac{8\pi^2}{15} \approx
        5.26$，此后递减。
  \item 证明 $V_n \to 0$ 比任何等比数列都快，
        并证明级数 $\sum_{n\geq1} V_n$ 收敛：所有
        单位球合起来的总体积有限。
  \item 证明母函数恒等式
        \[
        \sum_{k\geq0} V_{2k}\,x^{2k} = \eu^{\pi x^2}
        \qquad (x \in \R),
        \]
        并由此得出 $\sum_{k\geq0}V_{2k} = \eu^\pi \approx
        23.14$。
\end{enumerate}

\medskip
\noindent\textbf{第三部分 --- 第二条途径：$\Gamma$ 函数与
高斯积分。}
\begin{enumerate}[resume]
  \item 用富比尼定理（被积函数为乘积）证明
        \[
        I_n = \int_{\R^n}\eu^{-\norm
        x^2}\dd x
        = \Bigl(\int_{-\infty}^{+\infty}
        \eu^{-t^2}\dd t\Bigr)^{\!n} = \pi^{n/2},
        \]
        即 $n$ 维高斯积分，理解为
        在立方体 $\intcc{-R}{R}^n$ 上的极限。
  \item 回顾 $\Gamma(s) = \int_0^\infty t^{s-1}\eu^{-t}\dd
        t$（见\cref{def:b2:integration:gamma}）。
        由 $\Gamma(s+1) = s\,\Gamma(s)$
        （见\cref{thm:b2:integration:gammaprops}）及
        $\Gamma(\tfrac12) = \sqrt\pi$（作代换 $t = u^2$
        并利用高斯积分），计算
        \[
        \Gamma(k + 1) = k!,
        \qquad
        \Gamma\Bigl(k + \frac32\Bigr) =
        \frac{1\cdot3\cdots(2k+1)}{2^{k+1}}\,\sqrt\pi .
        \]
  \item 用归纳法，借助
        第 5 题的递推式，证明统一公式
        \[
        V_n = \frac{\pi^{n/2}}{\Gamma\bigl(\frac n2 +
        1\bigr)} \qquad (n \geq 1),
        \]
        并验证它能给出
        第 6 题中的两个闭式。
  \item 证明 $\int_0^\infty \eu^{-r^2}r^{n-1}\dd r =
        \tfrac12\Gamma\bigl(\tfrac n2\bigr)$，并由此推出
        恒等式
        \[
        I_n = n\,V_n\int_0^\infty \eu^{-r^2}\,r^{n-1}\,\dd r.
        \]
        解释如下：$\R^n$ 上的高斯质量
        沿球壳聚集，这些球壳的
        “$(n-1)$ 维面积”在半径 $r$ 处为
        $nV_nr^{n-1}$ --- 两边现已分别独立证明，
        因此这一解释无需额外论证。
  \item 记 $s_{n-1} = nV_n$（即单位球面
        $S^{n-1}$ 的面积，与 $v_n(R) = \int_0^R
        s_{n-1}r^{n-1}\dd r$ 相一致）。列出 $s_0, \dots, s_3$ 的值，并
        验证 $s_1 = 2\pi$，$s_2 = 4\pi$，$s_3 = 2\pi^2$。
\end{enumerate}

\medskip
\noindent\textbf{第四部分 --- 高维空间十分奇特。}
\begin{enumerate}[resume]
  \item 由斯特林公式
        （见\cref{thm:b2:comparison:stirling}）应用于 $k!$，
        证明对偶数 $n = 2k$，有
        \[
        V_n \sim \frac{1}{\sqrt{\pi n}}
        \Bigl(\frac{2\pi\eu}{n}\Bigr)^{n/2}
        \qquad (n \to \infty, \ n \text{ 偶数}),
        \]
        并解释为何同样的超等比衰减界
        借助递推式可推广到奇数 $n$。
  \item 单位球位于立方体 $\intcc{-1}1^n$ 中，
        此立方体的体积为 $2^n$。计算填充比 $V_n/2^n$ 在
        $n = 2, 3, 10$ 时的值，并证明它趋于 $0$：在高
        维情形下，立方体的几乎全部体积都集中在
        其角落中。
  \item 证明 $v_n(1)$ 中距边界球面距离
        在 $\varepsilon$ 以内的部分所占比例为 $1 -
        (1 - \varepsilon)^n \to 1$；数值上，
        $100$ 维球中落在
        厚度为 $1\%$ 的外壳内的比例是多少？
  \item 证明沃利斯渐近式 $W_n \sim
        \sqrt{\dfrac{\pi}{2n}}$ \emph{（$(W_n)$ 的单调性，比值
        $W_n/W_{n-2} \to 1$，以及 $W_nW_{n-1}
        = \frac\pi{2n}$）}，以及下界 $W_n \geq
        \sqrt{\dfrac{\pi}{2(n+1)}}$ 对所有 $n$ 成立。
  \item （平板上的集中）单位球中第一坐标超过
        $\delta$ 的部分所占比例为
        $\int_\delta^1(1 - x^2)^{\frac{n-1}2}\dd x \,\big/\,
        (2W_n)$。利用 $1 - u \leq \eu^{-u}$ 及尾部
        估计 $\int_\delta^\infty \eu^{-a x^2}\dd x \leq
        \frac{\eu^{-a\delta^2}}{2a\delta}$，证明此比例
        不超过
        \[
        \frac{\eu^{-(n-1)\delta^2/2}}{(n-1)\,\delta}
        \Big/ \sqrt{\frac{2\pi}{n+1}}
        \]
        由此得出结论：对 $\delta = s/\sqrt{n-1}$，
        除占比 $O(\eu^{-s^2/2}/s)$ 的部分外，球的其余部分都位于
        平板 $\abs{x_1} \leq s/\sqrt{n-1}$ 之内。高维球
        在统计意义上，同时在每个方向上都像
        一张薄饼。
  \item 将第 16--19 题综合为一段：说明
        $B_n(1)$ 的体积集中在何处（靠近边界
        球面，但仍落在过中心的每个超平面的 $O(1/\sqrt n)$
        厚平板之内），以及为何这两个
        结论互不矛盾。
\end{enumerate}

\medskip
\noindent\textbf{第 V 部分 --- 其他形体与综合。}
\begin{enumerate}[resume]
  \item （单纯形）设 $\Delta_n = \{x \in \R^n : x_i \geq 0,\
        \sum x_i \leq 1\}$。用切片与归纳法证明
        $\operatorname{vol}(\Delta_n) = \frac1{n!}$。
  \item （交叉多面体）由此可得 $C_n = \{x :
        \sum\abs{x_i} \leq 1\}$ 的体积为 $\frac{2^n}{n!}$，
        并从体积角度验证夹逼关系 $C_n \subseteq B_n(1)
        \subseteq \intcc{-1}1^n$，即
        $\frac{2^n}{n!} \leq V_n \leq 2^n$。
  \item 用第三种方法计算 $V_4$：将 $\R^4 = \R^2 \times
        \R^2$，对 $(z, w)$ 圆盘的面积在
        $(x, y)$ 圆盘上作极坐标积分，并
        还原出 $V_4 = \frac{\pi^2}2$。
  \item （蒙特卡罗方法失效）在
        立方体 $\intcc{-1}1^{20}$ 中均匀取一点，证明该点
        落入内切球的概率为
        $V_{20}/2^{20} \approx 2.5\cdot10^{-8}$，因此平均要
        约四千万次抽样才能首次命中，
        故用拒绝抽样估计 $V_n$
        在高维情形下失效（
        维数灾难）。
  \item 综合。两种互相独立的推导汇于 $V_n =
        \pi^{n/2}/\Gamma(\frac n2 + 1)$：分别列出两种推导所用的本章定理
        （富比尼定理、变量替换、
        极坐标下的高斯积分），并列出
        所用的一元结论（沃利斯公式、$\Gamma$ 函数、斯特林公式）。
        本科三年级分册在何处重做这一计算，
        借助勒贝格理论，它又增加了什么？
\end{enumerate}
\end{problem}
